% !TEX TS-program = pdflatex
% !TeX program = pdflatex
% !TEX encoding = UTF-8
% !TEX spellcheck = fr

\documentclass{beamer}
%\usepackage{fullpage}
%\usepackage[left=2.8cm,right=2.2cm,top=2 cm,bottom=2 cm]{geometry}
\setbeamersize{text margin left=10pt,text margin right=10pt}
\usepackage{amsmath,amssymb} 
\usepackage[T1]{fontenc}
\usepackage[utf8]{inputenc}
\usepackage[english,french]{babel}
\usepackage{txfonts}
\usepackage[]{graphicx}
\usepackage{multirow}
\usepackage{hyperref}

\hypersetup{
	colorlinks,
	urlcolor = blue
}

%\renewcommand{\baselinestretch}{1.5}

\def\supit#1{\raisebox{0.8ex}{\small\it #1}\hspace{0.05em}}

\title[BWEB: 01- Environnement] %
{Bureautique et Web \\Chapitre 01: Environnement de travail} 
\institute{ %
École  nationale Supérieure d'Informatique (ESI, ex. INI), Algérie
}
\author[ \textbf{\footnotesize  \insertframenumber/\inserttotalframenumber} \hspace*{\fill} ESI (2019-2020)] %
{ARIES Abdelkrime}
%\titlegraphic{\includegraphics[height=1cm]{../img/esi-logo.png}%\hspace*{4.75cm}~


\date{Année unniversitaire: 2019/2020} %\today

\usetheme{Warsaw} % Antibes Boadilla Warsaw

\beamertemplatenavigationsymbolsempty

%\setbeamertemplate{headline}{}


\begin{document}

\begin{frame}[plain]
	\maketitle
\end{frame}

\section{Introduction aux environnements de travail}

\begin{frame}
\frametitle{Introduction aux environnements de travail}

Un environnement de travail doit avoir:
\begin{itemize}
	\item un poste de travail: l'ordinateur et ses périphériques
	\item des applications: 
	\item espaces de stockage
\end{itemize}

\end{frame}

\subsection{Composition d'un environnement de travail}

\begin{frame}
\frametitle{Introduction aux environnements de travail}
\framesubtitle{Composition d'un environnement de travail}
%TODO schéma
\begin{itemize}
	\item Processeur 
	\item RAM 
	\item Stockage
	\item Périphériques 
	\item Système d'exploitation
	\item Applications 
\end{itemize}

\end{frame}

\subsection{Types d'environnements de travail}

\begin{frame}
\frametitle{Introduction aux environnements de travail}
\framesubtitle{Types d'environnements de travail}

Selon l'emplacement des données et des applications, un environnement de travail peut être: 
\begin{itemize}
	\item local
	\item distant
\end{itemize}

\end{frame}

\section{Environnement de travail local}

\begin{frame}
\frametitle{Environnement de travail local}

\end{frame}

\subsection{Matériel informatique}

\begin{frame}
\frametitle{Environnement de travail local}
\framesubtitle{Matériel informatique}

\end{frame}

\subsection{Système d'exploitation}

\begin{frame}
\frametitle{Environnement de travail local}
\framesubtitle{Système d'exploitation}

Parmi les fonctions d'un système d'exploitation: 
\begin{itemize}
	\item 
\end{itemize}

\end{frame}

\begin{frame}
\frametitle{Environnement de travail local}
\framesubtitle{Système d'exploitation}

SE de bureau (Desktop OS):
\begin{itemize}
	\item Windows (Microsoft)
	\item MacOS (Apple)
	\item Linux (Linux foundation)
	\item Chrome OS (Google): basé sur Linux
	\item BSD
\end{itemize}

SE de mobiles (Mobile OS): 
\begin{itemize}
	\item Android (Google)
	\item iOS (Apple)
	\item Windows 10 Mobile (Microsoft) 
	\item Chrome OS (Google)
	\item BSD
\end{itemize}

\end{frame}

\begin{frame}
\frametitle{Environnement de travail local}
\framesubtitle{Système d'exploitation}

%D'après \url{https://gs.statcounter.com/os-market-share/desktop/worldwide/#monthly-201901-201911}
% https://www.w3schools.com/browsers/browsers_os.asp

%\includegraphics[width=\textwidth]{imagefile}

\end{frame}

\begin{frame}
\frametitle{Environnement de travail local: Windows}
\framesubtitle{Système d'exploitation: Windows}

\end{frame}

\begin{frame}
\frametitle{Environnement de travail local: Windows}
\framesubtitle{Système d'exploitation: Windows}

\end{frame}

\begin{frame}
\frametitle{Environnement de travail local}
\framesubtitle{Système d'exploitation: Linux}

\end{frame}

\begin{frame}
\frametitle{Environnement de travail local: Linux}
\framesubtitle{Distributions}

Quelques distributions:

\begin{tabular}{llll}
	\includegraphics[width=.2\textwidth]{..//img/Bweb01-environnement/mint-logo.png} & \includegraphics[width=.2\textwidth]{..//img/Bweb01-environnement/ubuntu-logo.png} &
	\includegraphics[width=.2\textwidth]{..//img/Bweb01-environnement/opensuse-logo.png} &
	\includegraphics[width=.2\textwidth]{..//img/Bweb01-environnement/arch-logo.png} \\
	
	\includegraphics[width=.2\textwidth]{..//img/Bweb01-environnement/gentoo-logo.png} & 
	\includegraphics[width=.2\textwidth]{..//img/Bweb01-environnement/slackware-logo.png} &
	\includegraphics[width=.2\textwidth]{..//img/Bweb01-environnement/debian-logo.png} & 
	\includegraphics[width=.2\textwidth]{..//img/Bweb01-environnement/fedora-logo.png} \\
	
\end{tabular}

\end{frame}

\begin{frame}
\frametitle{Environnement de travail local: Linux}
\framesubtitle{Environnements de bureau}

\end{frame}

\begin{frame}
\frametitle{Environnement de travail local: Linux}
\framesubtitle{Environnements de bureau: KDE Plasma}

\begin{center}
	\includegraphics[height=
.8\textheight]{../img/Bweb01-environnement/kde.png}
\end{center}

\end{frame}

\begin{frame}
\frametitle{Environnement de travail local: Linux}
\framesubtitle{Environnements de bureau: GNOME3}

\begin{center}
	\includegraphics[height=
.8\textheight]{../img/Bweb01-environnement/gnome.png}
\end{center}

\end{frame}

\begin{frame}
\frametitle{Environnement de travail local: Linux}
\framesubtitle{Environnements de bureau: LXDE}

\begin{center}
	\includegraphics[height=
.8\textheight]{../img/Bweb01-environnement/lxde.png}
\end{center}

\end{frame}

\begin{frame}
\frametitle{Environnement de travail local: Linux}
\framesubtitle{Environnements de bureau: MATE}

\begin{center}
	\includegraphics[height=
.8\textheight]{../img/Bweb01-environnement/mate.png}
\end{center}

\end{frame}


\begin{frame}
\frametitle{Environnement de travail local: Linux}
\framesubtitle{Environnements de bureau: Cinnamon}

\begin{center}
	\includegraphics[height=
.8\textheight]{../img/Bweb01-environnement/cinnamon.png}
\end{center}

\end{frame}

\begin{frame}
\frametitle{Environnement de travail local: Linux}
\framesubtitle{Environnements de bureau: XFCE}

\begin{center}
	\includegraphics[height=
.8\textheight]{../img/Bweb01-environnement/xfce.png}
\end{center}

\end{frame}

\begin{frame}
\frametitle{Environnement de travail local}
\framesubtitle{Machine virtuelle (Virtual Box}

\begin{center}
	\includegraphics[height=
.8\textheight]{../img/Bweb01-environnement/vbox.png}
\end{center}

\end{frame}


\section{Environnement de travail distant}

\begin{frame}
\frametitle{Environnement de travail distant}

\end{frame}

\begin{frame}
\frametitle{Environnement de travail distant}

\end{frame}

\subsection{Accès à l'environnement de travail distant}

\begin{frame}
\frametitle{Environnement de travail distant}
\framesubtitle{Accès à l'environnement de travail distant}

\end{frame}

\begin{frame}
\frametitle{Environnement de travail distant: Accès}
\framesubtitle{Adressage (IP)}

Adresse IP

Local et Internet

\end{frame}

\begin{frame}
\frametitle{Environnement de travail distant: Accès}
\framesubtitle{Adressage (MAC)}

Adresse MAC

http://standards-oui.ieee.org/oui.txt



\end{frame}

\begin{frame}
\frametitle{Environnement de travail distant: Accès}
\framesubtitle{Adressage}

Comparaison IP vs. MAC



\end{frame}


\begin{frame}
\frametitle{Environnement de travail distant: Accès}
\framesubtitle{Site web}



\end{frame}


\begin{frame}
\frametitle{Environnement de travail distant: Accès}
\framesubtitle{URL (Uniform Resource Locator)}

Traduite: \textit{localisateur uniforme de ressource}, appelée \textit{adresse web}

{\huge \bfseries \textcolor{red}{https}://\textcolor{green}{www.w3schools.com}/browsers/browsers\_os.asp}

\begin{itemize}
	\item \textcolor{red}{protocole}
	\item \textcolor{green}{nom de domaine}
	\item 
\end{itemize}

\end{frame}


\begin{frame}
\frametitle{Environnement de travail distant: Accès}
\framesubtitle{Protocoles}

\begin{itemize}
	\item \textbf{HTTP (Hypertext Transfer Protocol)}: transfert hypertexte pour naviguer sur le web  (HTTPS pour la
	version sécurisée)
	\item \textbf{FTP (File Transfer Protocol)}: transfert des fichiers
	\item \textbf{SMTP (Simple Mail Transport Protocol)}: envoi et réception du courrier électronique
	\item \textbf{DNS (Domain Name System)}: transformation d'un URL vers une adresse IP
	\item \textbf{DHCP (Dynamic Host Configuration protocol)}: affectation des adresses IP à des clients (hôtes)
\end{itemize}

\end{frame}


\begin{frame}
\frametitle{Environnement de travail distant: Accès}
\framesubtitle{Serveurs}



\end{frame}

\begin{frame}
\frametitle{Environnement de travail distant: Accès}
\framesubtitle{Serveurs (Proxy}



\end{frame}


\begin{frame}
\frametitle{Environnement de travail distant: Accès}
\framesubtitle{Exercice: créer un serveur HTTP adhoc sous Windows}




\end{frame}


\subsection{Le cloud (le nuage)}

\begin{frame}
\frametitle{Environnement de travail distant}
\framesubtitle{Le cloud}

Le cloud est définit selon le NIST (National Institute of Standards and Technology)\footnote{La traduction en Français est fournie par \url{http://www.culture-informatique.net/cest-quoi-le-cloud/}}
\begin{definition}
	Le cloud computing est un modèle qui permet un accès omniprésent, pratique et à la demande à un réseau partagé et à un ensemble de ressources informatiques configurables (comme par exemple : des réseaux, des serveurs, du stockage, des applications et des services) qui peuvent être provisionnées et libérées avec un minimum d’administration ...
\end{definition}

\end{frame}

\begin{frame}
\frametitle{Environnement de travail distant: Le cloud}
\framesubtitle{Caractéristiques}

\begin{itemize}
	\item Un service à la demande
	\item Un accès aux ressources par le réseau
	\item Mise en commun des ressources
	\item Flexibilité des ressources
	\item Un service mesuré
\end{itemize}


\end{frame}

\begin{frame}
\frametitle{Environnement de travail distant: Le cloud}
\framesubtitle{Caractéristiques}

\begin{itemize}
	\item Un service à la demande
	\item Un accès aux ressources par le réseau
	\item Mise en commun des ressources
	\item Flexibilité des ressources
	\item Un service mesuré
\end{itemize}


\end{frame}


\begin{frame}
\frametitle{Environnement de travail distant: Le cloud}
\framesubtitle{Types}

\includegraphics[width=
\textwidth]{../img/Bweb01-environnement/cloud-types.png}


\end{frame}

\subsection{Google Drive}

\begin{frame}
\frametitle{Environnement de travail distant}
\framesubtitle{Google Drive}

\begin{itemize}
	\item un service de stockage et de partage de fichiers dans le cloud
	\item lancé par Google
	\item stockage gratuit de 15G (pour avoir plus, il faut payer)
	\item intègre une suite bureautique (Google Docs, Sheets, Slides, etc.)
	\item URL: \url{https://drive.google.com/}
\end{itemize}

\end{frame}


\begin{frame}
\frametitle{Environnement de travail distant: Google Drive}
\framesubtitle{Vue d'ensemble}

\begin{center}
	\includegraphics[width=
	\textwidth]{../img/Bweb01-environnement/drive.png}
\end{center}

\end{frame}

\begin{frame}
\frametitle{Environnement de travail distant: Google Drive}
\framesubtitle{Créer ou importer}

\begin{center}
	\includegraphics[height=
	.8\textheight]{../img/Bweb01-environnement/drive-new.png}
\end{center}

\end{frame}

\begin{frame}
\frametitle{Environnement de travail distant: Google Drive}
\framesubtitle{Paramètres d'un fichier (dossier)}

\begin{center}
	\includegraphics[height=
	.8\textheight]{../img/Bweb01-environnement/drive-options1.png}
	\includegraphics[height=
	.8\textheight]{../img/Bweb01-environnement/drive-options2.png}
\end{center}

\end{frame}

\begin{frame}
\frametitle{Environnement de travail distant: Google Drive}
\framesubtitle{Gestion des versions (fichiers no supportés par la suite Google) }

\begin{center}
	\includegraphics[height=
	.8\textheight]{../img/Bweb01-environnement/drive-version.png}
\end{center}

\end{frame}

\begin{frame}
\frametitle{Environnement de travail distant: Google Drive}
\framesubtitle{Gestion des versions (fichiers supportés par la suite Google) }

\begin{center}
	\includegraphics[width=
	.45\textwidth]{../img/Bweb01-environnement/drive-versionB1.png}
	\vline
	\includegraphics[width=
	.45\textwidth]{../img/Bweb01-environnement/drive-versionB2.png}
\end{center}

\end{frame}

\begin{frame}
\frametitle{Environnement de travail distant: Google Drive}
\framesubtitle{Partage de fichiers (dossiers) }

\begin{center}
	\includegraphics[width=
	.45\textwidth]{../img/Bweb01-environnement/drive-share1.png}
	\includegraphics[width=
	.45\textwidth]{../img/Bweb01-environnement/drive-share2.png}
\end{center}

\end{frame}


%\subsection{Bibliography}
%\frame[allowframebreaks]%
%{\frametitle{Bibliography}
%\tiny
%\bibliography{biblio}
%\bibliographystyle{apalike} 
%}


\end{document}

